\section{Conclusion}\label{sec:conclusion}

This paper shifts the evaluation of SME set-asides from access alone to
price formation. A set-aside does not merely direct contracts toward
protected firms; it changes the bidder pool from which the auction price
is formed. When the excluded non-SMEs are the bidders that discipline the
relevant order statistic, SME access is purchased by weakening the
competitive mechanism itself.

The paper does not estimate a full reduced-form treatment effect of the
legal reinterpretation, nor does it solve a dynamic entry model of SME
development. It identifies a narrower object: how the price-forming order
statistic changes when an exclusionary rule removes non-SMEs and the
protected SME pool responds. Within that static auction margin, exclusion
is expensive in standardized non-pharmaceutical procurement. A planner
can still prefer set-asides, but the justification must come from
benefits outside the measured static margin: SME surplus weights,
market-access objectives, or dynamic capacity gains.

The evidence from S\~ao Paulo's procurement reform shows that the
protected pool responds, but incompletely. SME participation rises after
the expansion of SME-only tendering, yet the post-policy protected pool
does not recreate the price discipline supplied by excluded non-SMEs. In
standardized non-pharmaceutical procurement, this leaves a large static
cost: the full set-aside generates a welfare loss of
\welfLossPctLthirtyNp{} percent of the open-regime price at
$\lambda=\lambdaMain$. Pharmaceutical procurement exhibits the boundary
of the exercise. There, the protected pool is thinner, composition
changes more, and the welfare ranking becomes model-sensitive.

The policy implication is that the instrument matters. A full set-aside
is a high-redistribution instrument because it forces eligibility by
removing rival bidders. A price preference delivers less redistribution,
but it preserves the non-SME bidders that discipline the price-forming
order statistic. The relevant comparison runs not between SME support
and no support, but between exclusionary redistribution and support that
keeps the price-forming bidder pool inside the auction. The evidence
does not imply that an implemented preference regime would reproduce the
static benchmark exactly; it shows that exclusion is the costly margin,
and that preserving the price-forming pool is valuable before any
behavioral response is modeled.

In thick standardized markets, where a preference rule can tilt
allocation without deleting rival bidders, exclusion is an expensive way
to redistribute. Support for SMEs need not remove the bidders that
discipline the price.

Two extensions would sharpen the result. The drop-out decomposition can be
run on other open reverse-auction platforms---federal ComprasNet is the
natural next case---to test whether exclusion dominance travels beyond a
single state and a single product group. And linking winning SMEs to firm
panels would let the dynamic capacity gains that sit outside the static
margin here be measured directly, turning the welfare-weight threshold into
an estimate rather than an indifference point.
